% !TeX root = ../main.tex
\chapter{非线性滤波与最优输运基础}[Foundations of Nonlinear Filtering and Optimal Transport]
% !TeX root = ../../main.tex
\section{符号定义与说明}\label{cw:app:notation}

\textbf{取值空间与密度导数：}随机变量 $X,Y$ 的取值空间分别记为 $\mathcal S_X,\mathcal S_Y$（上下文讨论单一空间且无歧义时简记为 $\mathcal S$），其上可测集合族记为 $\mathcal B(\mathcal S)$。有限维实向量空间上以体积测度为默认参考测度（状态空间上的体积测度记为 $v_X$），密度导数记为 $p_X(x)=\frac{\mathrm d\mathbb P_X}{\mathrm dv_X}(x)$，条件密度简写时记为 $p(x\mid y)=p_{X\mid Y}(x\mid y)$；一般测度间的密度导数记作 $\frac{\mathrm d\mu}{\mathrm d\nu}$，绝对连续记作 $\mu\ll\nu$。数学期望与概率带分布下标写作 $\mathbb E_{X\sim\mathbb P_X}[\cdot]$ 与 $\mathbb P_X(\cdot)$。

\textbf{概率单纯形与全一向量：}对正整数 $n$，$\mathbb R_+=[0,\infty)$，定义
\begin{equation*}
\Delta^{n-1}:=\left\{\bm w\in\mathbb R_+^n:\sum_{i=1}^n w_i=1\right\},\qquad
\mathbf1_n:=(1,\dots,1)^{\mathrm T}.
\end{equation*}
$\Delta^{n-1}$ 是 $n$ 个非负概率权重组成的集合，上标表示其维数；允许零权重，严格正权重另行说明。

\textbf{单点质量测度与示性函数：}对 $x\in\mathcal S$ 和可测集合 $A\in\mathcal B(\mathcal S)$，
\begin{equation*}
\delta_x(A):=\mathbb I_A(x)=
\begin{cases}1,&x\in A,\\0,&x\notin A,\end{cases}
\qquad \int_{\mathcal S}f(s)\,\mathrm d\delta_x(s)=f(x).
\end{equation*}
$\delta_x$ 是集中于单点 $x$ 的概率测度，与表示概率单纯形的大写 $\Delta$ 不同。

\textbf{概率测度与经验测度：}$\mathcal M(\mathcal S)$ 表示 $\mathcal S$ 上的全体非负测度，概率测度全体记为 $\mathcal M^1(\mathcal S):=\{\mu\in\mathcal M(\mathcal S):\mu(\mathcal S)=1\}$，不得将两者混用。记单点质量测度行数组 $\delta_{x_{1..N}}:=(\delta_{x_1},\dots,\delta_{x_N})$，对权重列向量 $\bm w\in\Delta^{N-1}$，加权与等权经验测度分别记为
\begin{equation*}
\delta_{x_{1..N}}\bm w:=\sum_{i=1}^Nw_i\delta_{x_i},\qquad
\hat{\mathbb P}_{N,\bm w;X}:=\delta_{x_{1..N}}\bm w,\qquad
\hat{\mathbb P}_{N;X}:=\frac1N\delta_{x_{1..N}}\mathbf1_N=\frac1N\sum_{i=1}^N\delta_{x_i}.
\end{equation*}
这里 $x_i$ 是样本实现值，行数组右乘列向量表示测度的有限凸组合。

\textbf{索引排列与滤波阶段：}用 $..$ 表示行排列数组 $x_{1..N}:=(x_1,\dots,x_N)$，用 $:$ 表示列拼接序列 $y_{1:k}:=(y_1,\dots,y_k)^{\mathrm T}$。对样本随机列向量 $X_i\in\mathbb R^{d\times1}$，记
\begin{equation*}
X_{1..N}:=(X_1,\dots,X_N)\in\mathbb R^{d\times N},\qquad
\mathbb P_{X_{1..N}}\text{ 为该样本族的联合分布}.
\end{equation*}
一般二元随机变量的联合分布记为 $\mathbb P_{(X,Y)}$。
上标 $-$ 与 $+$ 分别表示预测与后验阶段。位置变换中的 $X_k$ 表示按列拼接的粒子实现矩阵，每列是一个 $n_x$ 维粒子：
\begin{equation*}
\begin{aligned}
&X_k:=x_k^{(1..N)}:=[x_k^{(1)},\dots,x_k^{(N)}]\in\mathbb R^{n_x\times N},\\
&X_k^\pm:=x_k^{\pm,(1..N)}\in\mathbb R^{n_x\times N}.
\end{aligned}
\end{equation*}
$(\bm T_k)_{ij}$ 表示源粒子 $i$ 对目标粒子 $j$ 的贡献系数，故 $X_k^+=X_k^-\bm T_k$，$\bm T_k=N\bm\Pi_k$。概率分布下标中的 $X_k$ 仍表示状态随机变量，由上下文区分。

\textbf{推前、拉回与对偶配对：}对可测映射 $T:\mathcal S_X\to\mathcal S_Y$，
\begin{equation*}
\begin{aligned}
&(T_*\mu)(B):=\mu(T^{-1}(B)),\qquad T^*f:=f\circ T,\\
&\langle f,\mu\rangle_{\mathcal S}:=\int_{\mathcal S}f(s)\,\mathrm d\mu(s),\qquad
\langle f,T_*\mu\rangle_{\mathcal S_Y}=\langle T^*f,\mu\rangle_{\mathcal S_X}.
\end{aligned}
\end{equation*}
其中 $B\in\mathcal B(\mathcal S_Y)$，函数与测度的对偶配对尖括号须标明测度所在空间，积分对非负可测或相应可积函数使用。$\pi_i$ 表示取第 $i$ 个分量的坐标投影，$\mathbb P_X=X_*\mathbb P$。

\textbf{核复合与测度-核乘积耦合：}对条件转移核 $\kappa:\mathcal S_X\to\mathcal M(\mathcal S_Y)$ 与测度 $\mu\in\mathcal M(\mathcal S_X)$，复合测度记作 $\kappa\circ_m\mu:=\int_{\mathcal S_X}\kappa(x)\,\mu(\mathrm dx)\in\mathcal M(\mathcal S_Y)$；测度-核乘积耦合记作 $\mu\otimes_m\kappa\in\mathcal M(\mathcal S_X\times\mathcal S_Y)$，由此联合分布写作 $\mathbb P_{(X,Y)}=\mathbb P_X\otimes_m\mathbb P_{Y\mid X}$，独立乘积测度 $\mu\otimes\nu$（如 $\mathbb P_X\otimes\mathbb P_Y$）为其特例。

\textbf{联合概率矩阵与条件概率矩阵：}在正质量离散支撑上，记
\begin{equation*}
\begin{aligned}
&a_i:=\mathbb P_X(\{x_i\}),\qquad b_j:=\mathbb P_Y(\{y_j\}),\qquad
\Pi_{ij}:=\mathbb P_{(X,Y)}(\{(x_i,y_j)\}),\\
&P_{Y\mid X;i,j}:=\mathbb P_{Y\mid X=x_i}(\{y_j\})=\Pi_{ij}/a_i,\qquad
P_{X\mid Y;j,i}:=\mathbb P_{X\mid Y=y_j}(\{x_i\})=\Pi_{ij}/b_j.
\end{aligned}
\end{equation*}
条件矩阵的行对应条件变量，列对应被条件化变量；$\bm\Pi=(\Pi_{ij})$ 是联合概率矩阵。$\Pi(\mathbb P_X,\mathbb P_Y)$ 表示具有指定边缘的耦合测度集合，$U(\bm a,\bm b)$ 为其离散矩阵表示，见式~\eqref{cw:eq:coupling-family}、\eqref{cw:eq:transport-polytope}。

\textbf{矩阵与向量运算：}矩阵转置写作 $\bm A^{\mathrm T}$，迹写作 $\operatorname{tr}(\bm A)$，$\operatorname{diag}(\bm u)$ 将向量 $\bm u$ 置于矩阵对角线上；$\odot$、$\oslash$ 分别表示逐元素乘、除，$\otimes_k$ 表示矩阵分块直积。实矩阵迹内积记为 $\langle\bm A\mid\bm B\rangle:=\operatorname{tr}(\bm A^{\mathrm T}\bm B)=\sum_{i,j}A_{ij}B_{ij}$；向量一范数记为 $\|\bm z\|_1:=\sum_i|z_i|$，$\|\bm z\|$ 表示向量二范数。

\textbf{输运代价与熵正则化：}离散代价矩阵记为 $\bm C=(C_{ij})$，其中 $C_{ij}:=c(x_i,y_j)$。相对熵记为 $D_{\mathrm{KL}}(\mu\|\nu):=\int\log(\mathrm d\mu/\mathrm d\nu)\,\mathrm d\mu$（$\mu\ll\nu$ 时），否则为 $+\infty$；$\varepsilon>0$ 为熵正则化参数，$\tau>0$ 为边缘残差容差。

\textbf{平方可积空间与正交投影：}取值于 $\mathbb R^{n_x}$ 的平方可积随机变量空间记为 $\mathcal H:=L^2(\Omega,\mathcal F,\mathbb P;\mathbb R^{n_x})$；闭线性子空间关系记为 $\mathcal W\leqslant\mathcal H$，向 $\mathcal W$ 的正交投影算子与余投影算子分别记为 $\mathscr P_{\mathcal W}$ 与 $\mathscr P_{\mathcal W}^{\perp}$。

\textbf{参数空间、对偶伴随与压缩还原算子：}模型参数随机变量 $\bm\Theta$ 的取值空间记为 $\mathcal S_{\bm\Theta}$，参数实现记为 $\bm\theta\in\mathcal S_{\bm\Theta}$，含参条件密度记为 $p_{Y\mid(Z,\bm\Theta)}(y\mid z,\bm\theta)$。赋范或内积空间上的连续线性对偶与内积伴随算子统一记作上标 ${}^\star$，优化问题的最优解统一记作空心星上标 ${}^{\optstar}$（如 $\bm\theta^{\optstar},\mathscr T_{\Rightarrow}^{\optstar}$），以严格区分伴随运算与最优性。高维观测空间 $\mathcal S_X=\mathbb R^{n_D}$ 与低维潜空间 $\mathcal S_Z=\mathbb R^{n_d}$（$n_d<n_D$）之间的压缩算子与还原算子分别记为 $\mathscr T_{\Rightarrow}:\mathcal S_X\to\mathcal S_Z$ 与 $\mathscr T_{\Leftarrow}:\mathcal S_Z\to\mathcal S_X$，线性情形下二者在标准基下的表示矩阵记为 $\bm T_{\Rightarrow},\bm T_{\Leftarrow}$。


% !TeX root = ../../main.tex
\section{非线性滤波问题}\label{cw:sec:model}

本文从状态空间模型与 Bayes 递推出发，梳理粒子滤波及最优输运重采样的数学关系，并给出算法实现与仿真比较。经典滤波框架与序贯蒙特卡罗方法分别参考文献\cite{bar2001estimation,doucet2001sequential}。

本文约定：凡使用密度公式时，均假设相应密度存在，以相应空间的体积测度为参考测度，密度等式按几乎处处理解。状态空间的体积测度记为 $v_X$。第~\ref{cw:sec:sim-logistic} 节的零过程噪声模型使用单点质量转移核，预测按测度推前理解。

设概率空间为 $(\Omega,\mathcal{F},\mathbb{P})$，状态空间为 $(\mathcal{S}_X,\mathcal{B}(\mathcal{S}_X))$，观测空间为 $(\mathcal{S}_Y,\mathcal{B}(\mathcal{S}_Y))$，其中 $\mathcal{S}_X=\mathbb{R}^{n_x}$，$\mathcal{S}_Y=\mathbb{R}^{n_y}$，$\mathcal{B}$ 表示可测集合族。状态与观测满足
\begin{equation}\label{cw:eq:1}
\begin{aligned}
&X_k=f(X_{k-1},N_{x,k}),\\
&Y_k=h(X_k,N_{y,k}),\;k\geq1.
\end{aligned}
\end{equation}
其中 $X_k:\Omega\to\mathcal{S}_X$，$Y_k:\Omega\to\mathcal{S}_Y$ 为可测映射，$f,h$ 为已知可测函数。$\mathbb{P}$ 定义在 $\Omega$ 上，状态分布 $\mathbb{P}_{X_k}=(X_k)_*\mathbb{P}$ 与滤波后验分布定义在 $\mathcal{S}_X$ 上。

条件：
\begin{itemize}
\item 所有时刻的过程噪声与观测噪声彼此独立，并与 $X_0$ 独立，即 $X_0, N_{x,k}, N_{y,k}$ 作为随机变量相互独立。
\item 观测条件密度 $p_{Y_k\mid X_k}(y \mid x)$ 可计算，对实际观测，Bayes 更新的归一化常数有限且严格为正。
\end{itemize}

输入：模型 $f,h$，初始分布 $\mathbb{P}_{X_0}$，噪声分布 $\mathbb{P}_{N_x},\mathbb{P}_{N_y}$，依次到达的观测值 $y_1,y_2,\dots$。

输出：逐个收到观测 $y_k$，递推求解当前状态的后验概率密度函数 $p_{X_k\mid Y_{1:k}}(x_k \mid y_{1:k})$。

% !TeX root = ../../main.tex
\section{Bayes 滤波}\label{cw:sec:bayes}

\begin{lemma}[无记忆性\cite{sarkka2013bayesian,cappe2005inference}]\label{cw:lem:markov}
在上述状态方程与噪声独立性条件下，给定 $X_{k-1}$ 后，$X_k$ 的条件分布不再依赖更早的状态与历史观测；给定 $X_k$ 后，$Y_k$ 的条件分布不再依赖过去的状态与观测。对 $k\geq1$ 有
\begin{equation}\label{cw:eq:2}
\begin{aligned}
&p_{X_k\mid X_{0:k-1},Y_{1:k-1}}(x_k\mid x_{0:k-1},y_{1:k-1})
=p_{X_k\mid X_{k-1}}(x_k\mid x_{k-1}),\\
&p_{Y_k\mid X_{0:k},Y_{1:k-1}}(y_k\mid x_{0:k},y_{1:k-1})
=p_{Y_k\mid X_k}(y_k\mid x_k).
\end{aligned}
\end{equation}
$k=1$ 时历史观测为空。第一式给出状态的单步无记忆性，第二式给出观测的条件独立性，分别用于下述预测与更新。

\end{lemma}
\begin{proof}
为完整起见，本文给出自行推导的证明：离散情形的证明及一般情形的推广说明见附录~\ref{cw:app:markov-proof}~节。
\end{proof}

固定观测记录 $y_{1:k}=(y_1,\dots,y_k)^{\mathrm{T}}$，以 $p_{X_0}$ 初始化，$k=1$ 时不含历史观测条件。

\textbf{预测：}由上一时刻的后验与状态转移密度，求当前状态在收到 $y_k$ 前的条件密度。
\begin{equation}\label{cw:eq:3}
\begin{aligned}
&p_{X_k\mid Y_{1:k-1}}(x_k\mid y_{1:k-1})\\
= &\int_{\mathcal{S}_X}p_{X_k\mid X_{k-1}}(x_k\mid x_{k-1})p_{X_{k-1}\mid Y_{1:k-1}}(x_{k-1}\mid y_{1:k-1})\,\mathrm{d}v_X(x_{k-1}).
\end{aligned}
\end{equation}

\textbf{更新：}收到 $y_k$ 后，用观测似然对预测密度加权，再归一化，证明见定理~\ref{cw:thm:bayes}。
\begin{equation}\label{cw:eq:4}
\begin{aligned}
&p_{X_k\mid Y_{1:k}}(x_k\mid y_{1:k})\\
= &\frac{p_{Y_k\mid X_k}(y_k\mid x_k)p_{X_k\mid Y_{1:k-1}}(x_k\mid y_{1:k-1})}
{\displaystyle\int_{\mathcal{S}_X}p_{Y_k\mid X_k}(y_k\mid x)p_{X_k\mid Y_{1:k-1}}(x\mid y_{1:k-1})\,\mathrm{d}v_X(x)}.
\end{aligned}
\end{equation}
分母就是观测的预测密度 $p_{Y_k\mid Y_{1:k-1}}(y_k\mid y_{1:k-1})$，要求其有限且严格为正。更新所得后验作为下一步预测的输入，这两步构成精确的 Bayes 滤波递推。

% !TeX root = ../../main.tex
\section{粒子滤波与经验测度近似}\label{cw:sec:particle}

以下以 Gordon 等提出的自举粒子滤波为基础\cite{gordon1993novel}，采用序贯蒙特卡罗的经验测度表示\cite{doucet2001sequential}。粒子滤波的理论框架、主要算法与多目标跟踪应用已有系统的中文综述\cite{li2015particlezh}。

固定观测记录，用 $-$ 与 $+$ 分别表示预测与后验，记
\begin{equation}\label{cw:eq:5}
\begin{aligned}
&\mathbb{P}_{X_k}^-:=\mathbb{P}_{X_k\mid Y_{1:k-1}=y_{1:k-1}},\\
&\mathbb{P}_{X_k}^+:=\mathbb{P}_{X_k\mid Y_{1:k}=y_{1:k}},\\
&\mathbb{P}_{X_0}^+:=\mathbb{P}_{X_0}.
\end{aligned}
\end{equation}

从预测公式出发，以下使用条件密度的简写 $p(x_k\mid y_{1:k-1})=p_{X_k\mid Y_{1:k-1}}(x_k\mid y_{1:k-1})$，其他条件密度同理：
\begin{equation}\label{cw:eq:6}
\begin{aligned}
&p(x_k\mid y_{1:k-1})\\
= &\int_{\mathcal{S}_X}p(x_k\mid x_{k-1})p(x_{k-1}\mid y_{1:k-1})\,\mathrm{d}v_X(x_{k-1}).
\end{aligned}
\end{equation}

固定 $x_k$，这个积分就是函数 $x_{k-1}\mapsto p(x_k\mid x_{k-1})$ 关于上一时刻后验分布的平均。因此，先假设能够从 $p(x_{k-1}\mid y_{1:k-1})$ 独立抽取 $N$ 个样本，记实现值为 $x_{k-1}^{+,(1)},\dots,x_{k-1}^{+,(N)}$，其中 $+$ 表示这些样本来自后验分布。

由简单抽样得到等权经验测度，每个样本占 $1/N$ 的质量：
\begin{equation}\label{cw:eq:7}
\hat{\mathbb{P}}_{N;X_{k-1}}^+
:=\frac1N\sum_{i=1}^N\delta_{x_{k-1}^{+,(i)}}.
\end{equation}
用这一经验测度替换预测积分中的后验测度 $p(x_{k-1}\mid y_{1:k-1})\,\mathrm{d}v_X(x_{k-1})$，得到
\begin{equation}\label{cw:eq:8}
\begin{aligned}
&p(x_k\mid y_{1:k-1})\\
\approx &\int_{\mathcal{S}_X}p(x_k\mid x_{k-1})\,\mathrm{d}\hat{\mathbb{P}}_{N;X_{k-1}}^+(x_{k-1})\\
= &\frac1N\sum_{i=1}^N\int_{\mathcal{S}_X}p(x_k\mid x_{k-1})\,\mathrm{d}\delta_{x_{k-1}^{+,(i)}}(x_{k-1}).
\end{aligned}
\end{equation}
最后使用单点质量测度的积分性质 $\int\varphi(x)\,\mathrm{d}\delta_a(x)=\varphi(a)$，每个积分都变成在对应样本处的函数值：
\begin{equation}\label{cw:eq:9}
p(x_k\mid y_{1:k-1})\approx\frac1N\sum_{i=1}^N p(x_k\mid x_{k-1}^{+,(i)}).
\end{equation}
这样，后验分布下的积分被样本平均近似，预测密度成为 $N$ 个转移密度的等权混合。

\textbf{提议分布：}在生成新粒子时，需要指定实际用于抽样的分布，其密度称为提议密度，记为 $q$。它可以依赖旧粒子与当前观测，抽样操作写作
\begin{equation}\label{cw:eq:10}
x_k^{-,(i)}\sim q(\cdot\mid x_{k-1}^{+,(i)},y_k).
\end{equation}
当前要对上面的预测混合分布进行采样，我们选择每个成分本身的转移密度作为提议密度：
\begin{equation}\label{cw:eq:11}
q(x_k\mid x_{k-1}^{+,(i)},y_k)=p(x_k\mid x_{k-1}^{+,(i)}).
\end{equation}
这一选择在生成粒子时不使用当前观测 $y_k$，正是自举滤波采用的转移先验提议。

\textbf{分层采样：}以混合成分的编号 $i$ 作为分层依据，第 $i$ 层对应密度 $p(x_k\mid x_{k-1}^{+,(i)})$，在混合分布中占 $1/N$ 的质量。因此每层分配一个样本，给定上一时刻的粒子后，各层独立抽取预测粒子：
\begin{equation}\label{cw:eq:12}
x_k^{-,(i)}\sim p(\cdot\mid x_{k-1}^{+,(i)}),\;i=1,\dots,N.
\end{equation}
这里 $\sim$ 表示抽样操作，$x_k^{-,(i)}$ 记抽得的实现值。由状态方程，这一操作可以通过独立抽取过程噪声实现值 $n_{x,k}^{(i)}$ 完成：
\begin{equation}\label{cw:eq:13}
x_k^{-,(i)}=f\bigl(x_{k-1}^{+,(i)},n_{x,k}^{(i)}\bigr).
\end{equation}

在上述转移先验提议下，每个样本代表所在层的 $1/N$ 质量，因此将这些预测粒子合并，得到预测分布的等权经验测度：
\begin{equation}\label{cw:eq:14}
\hat{\mathbb{P}}_{N;X_k}^-:=\frac1N\sum_{i=1}^N\delta_{x_k^{-,(i)}}
\approx\mathbb{P}_{X_k}^-.
\end{equation}
这一步用有限个点质量进一步近似前面得到的混合分布。给定旧粒子后，各层样本独立，但各自服从对应的转移分布。

\textbf{更新：}$-$ 与 $+$ 分别表示使用当前观测前后的预测与后验阶段，此处只更新权重。
对任意可测集合 $A$，将 Bayes 公式中的预测测度替换为等权经验测度，再用单点质量测度的积分性质展开；假设样本似然总和有限且严格为正，约去 $1/N$ 后得到加权后验近似：
\begin{equation}\label{cw:eq:particle-update}
\begin{aligned}
&\mathbb{P}_{X_k}^+(A)\\
={}&\int_A\frac{p(y_k\mid x)p(x\mid y_{1:k-1})}{\displaystyle\int_{\mathcal{S}_X}p(y_k\mid z)p(z\mid y_{1:k-1})\,\mathrm{d}v_X(z)}\,\mathrm{d}v_X(x)\\
={}&\frac{\displaystyle\int_A p(y_k\mid x)\,\mathrm{d}\mathbb{P}_{X_k}^-(x)}{\displaystyle\int_{\mathcal{S}_X}p(y_k\mid x)\,\mathrm{d}\mathbb{P}_{X_k}^-(x)},\\
\Rightarrow\quad &\hat{\mathbb{P}}_{N,\bm w_k^+;X_k}^+(A):=\frac{\displaystyle\int_A p(y_k\mid x)\,\mathrm{d}\hat{\mathbb{P}}_{N;X_k}^-(x)}{\displaystyle\int_{\mathcal{S}_X}p(y_k\mid x)\,\mathrm{d}\hat{\mathbb{P}}_{N;X_k}^-(x)}\\
={}&\sum_{i=1}^N w_k^{+,(i)}\delta_{x_k^{-,(i)}}(A),\\
&w_k^{+,(i)}:=\frac{p(y_k\mid x_k^{-,(i)})}{\displaystyle\sum_{j=1}^N p(y_k\mid x_k^{-,(j)})},\;i=1..N.
\end{aligned}
\end{equation}

\textbf{重采样：}
粗略地说，重采样是对粒子位置的更新。因此，和上一步的概率测度更新应该联为一体，同属于预测 $\to $ 更新的大步骤。
粒子矩阵按列排列为 $X_k^\pm:=x_k^{\pm,(1..N)}\in\mathbb R^{n_x\times N}$。用 $\bm\Pi_k$ 表示联合概率矩阵，$\bm T_k$ 表示位置变换矩阵，$\bm n_k$ 表示各粒子的复制次数，$(\bm T_k)_{ij}$ 表示源粒子 $i$ 对目标粒子 $j$ 的贡献系数。经典重采样中该系数为 $\mathbb I_{\{A_j=i\}}$，其中 $A_j$ 是目标粒子 $j$ 抽取的源索引：
\begin{equation}\label{cw:eq:resampling}
\begin{aligned}
&X_k^+=X_k^-\bm T_k=NX_k^-\bm\Pi_k,\;\bm T_k=N\bm\Pi_k\geq0,\\
&\bm T_k^{\mathrm T}\mathbf1_N=\mathbf1_N,\;\bm T_k\mathbf1_N=\bm n_k,\;\mathbb E[\bm n_k\mid\bm w_k^+]=N\bm w_k^+.
\end{aligned}
\end{equation}
重采样就是按后验权重，从已有预测粒子中有放回地抽取 $N$ 次。先把各粒子的权重依次累加，将 $[0,1)$ 划成 $N$ 段，每段长度就是对应粒子的权重。
每次独立抽一个均匀随机数，落在哪一段，就复制哪一个粒子；重复 $N$ 次，得到后验粒子。因而，权重越大的粒子越容易被多次保留，权重小的粒子可能一次也选不到。
将新粒子的权重统一设为 $1/N$，原来由权重表达的质量差异，就改由复制次数表达。
这样，$\bm T_k$ 的第 $j$ 列记录新粒子 $j$ 复制了谁，$\bm n_k$ 记录每个旧粒子被复制了几次。单次重采样会产生随机误差；固定原粒子与权重后，对重采样随机性取平均，所得经验测度等于原加权测度。

% !TeX root = ../../main.tex
\section{最优输运与重采样}\label{cw:sec:ot}

经典重采样通过随机复制选择耦合矩阵，其源边缘为实际复制比例，只有对重采样随机性取平均后才恢复后验权重：
\begin{equation}\label{cw:eq:ot-random-marginal}
\bm\Pi_k\mathbf1_N=\frac{\bm n_k}{N},\qquad
\bm\Pi_k^{\mathrm T}\mathbf1_N=\frac1N\mathbf1_N,\qquad
\mathbb E_{\bm n_k\sim\mathbb P_{\bm n_k\mid\bm w_k^+}}
\!\left[\frac{\bm n_k}{N}\right]=\bm w_k^+.
\end{equation}
另一种选择是直接约束源边缘为 $\bm w_k^+$，并按粒子间的移动代价选择耦合。为此，先从输运映射出发，建立联合分布、条件概率与离散耦合矩阵之间的关系。

\subsection{从 Monge 映射到 Kantorovich 耦合}\label{cw:sec:monge}

本节采用最优输运的标准映射与耦合表述\cite{villani2009optimal,peyre2019computational}；其基本理论、离散模型与正则化方法亦可参考文献\cite{ma2019optimaltransportzh}。

本节暂用 $X,Y$ 表示一般输运问题的源、目标随机变量，$Y$ 不特指滤波观测。给定边缘分布 $\mathbb P_X,\mathbb P_Y$ 与非负代价函数 $c$，Monge 问题寻找可测映射 $T$，使
\begin{equation}\label{cw:eq:monge}
\begin{aligned}
&Y=T(X),\qquad T_*\mathbb P_X=T_*X_*\mathbb P=\mathbb P_Y,\\
&\inf_{T:\,T_*\mathbb P_X=\mathbb P_Y}
\int_{\mathcal S_X}c(x,T(x))\,\mathrm d\mathbb P_X(x)
=\inf_{T:\,T_*\mathbb P_X=\mathbb P_Y}
\mathbb E_{X\sim\mathbb P_X}[c(X,T(X))].
\end{aligned}
\end{equation}
确定性绑定 $Y=T(X)$ 使条件分布退化为 $\mathbb P_{Y\mid X=x}=\delta_{T(x)}$：一个源点的质量只能整体送往一个目标点，不能拆分。离散边缘因而可能不存在这样的映射；均匀离散分布的整除条件见附录~\ref{cw:app:monge}。

\textbf{耦合松弛：}解除函数绑定，允许任意满足两个边缘约束的联合分布。记坐标投影 $\pi_1(x,y)=x$、$\pi_2(x,y)=y$，定义
\begin{equation}\label{cw:eq:coupling-family}
\Pi(\mathbb P_X,\mathbb P_Y)
:=\left\{\mu\in\mathcal M^1(\mathcal S_X\times\mathcal S_Y):
(\pi_1)_*\mu=\mathbb P_X,\ (\pi_2)_*\mu=\mathbb P_Y\right\}.
\end{equation}
输运目标于是成为联合分布上的期望最小化：
\begin{equation}\label{cw:eq:kantorovich}
\inf_{\mu\in\Pi(\mathbb P_X,\mathbb P_Y)}
\int_{\mathcal S_X\times\mathcal S_Y}c(x,y)\,\mathrm d\mu(x,y)
=\inf_{\mathbb P_{(X,Y)}\in\Pi(\mathbb P_X,\mathbb P_Y)}
\mathbb E_{(X,Y)\sim\mathbb P_{(X,Y)}}[c(X,Y)].
\end{equation}
在有限维实向量空间上，联合分布可通过条件转移核写成测度-核乘积耦合
$\mathbb P_{(X,Y)}=\mathbb P_X\otimes_m\mathbb P_{Y\mid X}$；独立乘积
$\mathbb P_X\otimes\mathbb P_Y$ 是其中一个特例。

\subsection{离散耦合与条件概率及边缘约束}\label{cw:sec:discrete-ot}

设两组支撑点各自互异，$\bm a\in\Delta^{m-1}$、$\bm b\in\Delta^{n-1}$，且暂取 $a_i,b_j>0$；符号定义见第~\ref{cw:app:notation}~节。记
\begin{equation}\label{cw:eq:discrete-coupling}
\begin{aligned}
&\mathbb P_X=\delta_{x_{1..m}}\bm a,\qquad
\mathbb P_Y=\delta_{y_{1..n}}\bm b,\\
&\Pi_{ij}:=\mathbb P_{(X,Y)}(\{(x_i,y_j)\}),\qquad
\mathbb P_{(X,Y)}=\sum_{i=1}^m\sum_{j=1}^n\Pi_{ij}\,
\delta_{x_i}\otimes\delta_{y_j}.
\end{aligned}
\end{equation}
这里 $\otimes$ 表示测度的乘积。对点质量作坐标推前，就得到边缘的行和、列和：
\begin{equation}\label{cw:eq:discrete-marginals}
\begin{aligned}
&(\pi_1)_*(\delta_{x_i}\otimes\delta_{y_j})=\delta_{x_i},\qquad
(\pi_2)_*(\delta_{x_i}\otimes\delta_{y_j})=\delta_{y_j},\\
&(\pi_1)_*\mathbb P_{(X,Y)}=\sum_{i=1}^m\left(\sum_{j=1}^n\Pi_{ij}\right)\delta_{x_i}=\mathbb P_X,\\
&(\pi_2)_*\mathbb P_{(X,Y)}=\sum_{j=1}^n\left(\sum_{i=1}^m\Pi_{ij}\right)\delta_{y_j}=\mathbb P_Y.
\end{aligned}
\end{equation}
函数与测度的对偶配对记为 $\langle f,\mu\rangle_{\mathcal S}:=\int_{\mathcal S}f(s)\,\mathrm d\mu(s)$，下标标明测度所在空间。
对于一般情况（假使需要讨论连续而非离散），等价地，用推前与拉回对偶，对任意可测集合 $A\in\mathcal B(\mathcal S_X)$ 有
\begin{equation}\label{cw:eq:marginal-duality}
\begin{aligned}
\langle\mathbb I_A\otimes\mathbf1_{\mathcal S_Y},\mathbb P_{(X,Y)}\rangle_{\mathcal S_X\times\mathcal S_Y}
&=\langle\pi_1^*\mathbb I_A,\mathbb P_{(X,Y)}\rangle_{\mathcal S_X\times\mathcal S_Y}\\
&=\langle\mathbb I_A,(\pi_1)_*\mathbb P_{(X,Y)}\rangle_{\mathcal S_X}
=\langle\mathbb I_A,\mathbb P_X\rangle_{\mathcal S_X}.
\end{aligned}
\end{equation}
\begingroup\raggedright
更一般地，每个坐标分别满足 $(\pi_i)_*\mathbb P_{X_{1..N}}=\mathbb P_{X_i}$；全部边缘组成的族应写为
$((\pi_i)_*\mathbb P_{X_{1..N}})_{i=1}^N=(\mathbb P_{X_i})_{i=1}^N$。\par
\endgroup

\textbf{条件概率：}条件概率矩阵标明条件方向，联合概率矩阵统一记为 $\bm\Pi$：
\begin{equation}\label{cw:eq:conditional-matrices}
\begin{aligned}
&P_{Y\mid X;i,j}:=\mathbb P_{Y\mid X=x_i}(\{y_j\})=\frac{\Pi_{ij}}{a_i},\qquad
P_{X\mid Y;j,i}:=\mathbb P_{X\mid Y=y_j}(\{x_i\})=\frac{\Pi_{ij}}{b_j},\\
&\bm\Pi=\operatorname{diag}(\bm a)\bm P_{Y\mid X}
=\bm P_{X\mid Y}^{\mathrm T}\operatorname{diag}(\bm b),\\
&\bm P_{Y\mid X}\mathbf1_n=\mathbf1_m,\qquad
\bm P_{Y\mid X}^{\mathrm T}\bm a=\bm b.
\end{aligned}
\end{equation}
条件矩阵以条件变量的取值编号为行、被条件化变量的取值编号为列。

\textbf{离散输运：}令 $C_{ij}=c(x_i,y_j)$，连续耦合约束化为输运多面体，目标化为有限求和：
\begin{equation}\label{cw:eq:transport-polytope}
\begin{aligned}
&U(\bm a,\bm b):=\{\bm\Pi\in\mathbb R_+^{m\times n}:
\bm\Pi\mathbf1_n=\bm a,\ \bm\Pi^{\mathrm T}\mathbf1_m=\bm b\},\\
&\min_{\bm\Pi\in U(\bm a,\bm b)}\langle\bm\Pi\mid\bm C\rangle,
\qquad\langle\bm\Pi\mid\bm C\rangle:=\operatorname{tr}(\bm\Pi^{\mathrm T}\bm C)=\sum_{i=1}^m\sum_{j=1}^n\Pi_{ij}C_{ij}.
\end{aligned}
\end{equation}
独立联合分布对应 $\bm\Pi=\bm a\bm b^{\mathrm T}\in U(\bm a,\bm b)$。

\subsection{熵正则化与 Sinkhorn 迭代}\label{cw:sec:sinkhorn}

采用 Cuturi 的熵正则化离散输运方案\cite{cuturi2013sinkhorn}，推导及数值实现参考文献\cite{peyre2019computational}。

\textbf{正则化：}给定 $\varepsilon>0$，以独立联合分布为参考测度，考虑
\begin{equation}\label{cw:eq:regularized-measure}
\mathbb P_{(X,Y)}^\varepsilon
=\operatorname*{argmin}_{\mu\in\Pi(\mathbb P_X,\mathbb P_Y)}
\left\{\int_{\mathcal S_X\times\mathcal S_Y} c(x,y)\,\mathrm d\mu(x,y)
+\varepsilon D_{\mathrm{KL}}(\mu\|\mathbb P_X\otimes\mathbb P_Y)\right\}.
\end{equation}
在离散边缘固定时，KL 项与负熵仅相差一个常数；以下采用 $0\log0=0$：
\begin{equation}\label{cw:eq:kl-entropy}
\begin{aligned}
D_{\mathrm{KL}}(\bm\Pi\|\bm a\bm b^{\mathrm T})
&=\sum_{i=1}^m\sum_{j=1}^n\Pi_{ij}\log\frac{\Pi_{ij}}{a_i b_j}\\
&=\sum_{i=1}^m\sum_{j=1}^n\Pi_{ij}\log\Pi_{ij}
-\sum_{i=1}^m a_i\log a_i-\sum_{j=1}^n b_j\log b_j.
\end{aligned}
\end{equation}
因此，所求矩阵等价地满足
\begin{equation}\label{cw:eq:regularized-matrix}
\bm\Pi^\varepsilon
=\operatorname*{argmin}_{\bm\Pi\in U(\bm a,\bm b)}
\left\{\langle\bm\Pi\mid\bm C\rangle
+\varepsilon\sum_{i=1}^m\sum_{j=1}^n\Pi_{ij}(\log\Pi_{ij}-1)\right\}.
\end{equation}
这里额外减去 $\varepsilon\sum_{i,j}\Pi_{ij}=\varepsilon$ 也是常数，使求导后的表达更简洁。

\textbf{对偶条件：}取 $\bm\alpha\in\mathbb R^m$、$\bm\beta\in\mathbb R^n$，构造拉格朗日函数：
\begin{equation}\label{cw:eq:sinkhorn-lagrangian}
\begin{aligned}
\mathcal L(\bm\Pi,\bm\alpha,\bm\beta)
={}&\langle\bm\Pi\mid\bm C\rangle
+\varepsilon\sum_{i=1}^m\sum_{j=1}^n\Pi_{ij}(\log\Pi_{ij}-1)\\
&+\bm\alpha^{\mathrm T}(\bm\Pi\mathbf1_n-\bm a)
+\bm\beta^{\mathrm T}(\bm\Pi^{\mathrm T}\mathbf1_m-\bm b).
\end{aligned}
\end{equation}
有限代价与严格正边缘下，最优矩阵各元素严格为正；一阶条件给出
\begin{equation}\label{cw:eq:sinkhorn-scaling}
\begin{aligned}
&\frac{\partial\mathcal L}{\partial\Pi_{ij}}
=C_{ij}+\varepsilon\log\Pi_{ij}+\alpha_i+\beta_j=0,\\
&K_{ij}:=\exp(-C_{ij}/\varepsilon),\qquad
u_i:=\exp(-\alpha_i/\varepsilon),\quad v_j:=\exp(-\beta_j/\varepsilon),\\
&\Pi_{ij}^\varepsilon=u_iK_{ij}v_j,\qquad
\bm\Pi^\varepsilon=\operatorname{diag}(\bm u)\bm K\operatorname{diag}(\bm v).
\end{aligned}
\end{equation}

\textbf{交替缩放：}将矩阵分解代回边缘约束，得到
\begin{equation}\label{cw:eq:sinkhorn-iteration}
\begin{aligned}
&\bm u\odot(\bm K\bm v)=\bm a,\qquad
\bm v\odot(\bm K^{\mathrm T}\bm u)=\bm b,\\
&\bm v^{(0)}=\mathbf1_n,\qquad
\bm u^{(\ell+1)}=\bm a\oslash(\bm K\bm v^{(\ell)}),\\
&\bm v^{(\ell+1)}=\bm b\oslash(\bm K^{\mathrm T}\bm u^{(\ell+1)}),\\
&\bm\Pi^{(\ell+1)}
=\operatorname{diag}(\bm u^{(\ell+1)})\bm K\operatorname{diag}(\bm v^{(\ell+1)}).
\end{aligned}
\end{equation}
其中 $\odot,\oslash$ 分别表示逐元素乘、除。每次先校正行边缘，再校正列边缘；以两侧边缘残差共同判断停止：
\begin{equation}\label{cw:eq:sinkhorn-stop}
\max\!\left\{\|\bm\Pi^{(\ell)}\mathbf1_n-\bm a\|_1,
\|(\bm\Pi^{(\ell)})^{\mathrm T}\mathbf1_m-\bm b\|_1\right\}\le\tau.
\end{equation}
严格正核下，迭代趋向唯一的正则化最优耦合。

\subsection{基于条件均值的粒子位置变换}\label{cw:sec:ot-resampling}

该构造沿用 Reich 的集合变换方法\cite{reich2013nonparametric}，并用熵正则化耦合近似未正则化的输运解。

沿用离散输运中的边际记号 $\bm a,\bm b$, 当前时刻的代价矩阵记为 $\bm C_k\in\mathbb R^{N\times N}$, 令
\begin{equation}\label{cw:eq:particle-ot-input}
\begin{aligned}
&\bm a=\bm w_k^+,\;\bm b=\frac1N\mathbf1_N,\;(\bm C_k)_{ij}:=\|x_k^{-,(i)}-x_k^{-,(j)}\|^2,\\
&\bm\Pi_k\in U(\bm a,\bm b),\;\bm\Pi_k\mathbf1_N=\bm w_k^+,\;\bm\Pi_k^{\mathrm T}\mathbf1_N=\frac1N\mathbf1_N.
\end{aligned}
\end{equation}
这里 $i,j=1,\dots,N$, 第 $i$ 行对应后验权重为 $w_k^{+,(i)}$ 的预测粒子, 第 $j$ 列对应等权的目标参考位置 $x_k^{-,(j)}$。源边际为加权后验经验测度, 目标边际为等权预测经验测度。将 $\bm a,\bm b,\bm C_k$ 代入 式~\eqref{cw:eq:regularized-matrix} 的熵正则化目标, 用 Sinkhorn 求解 $\bm\Pi_k$；正则化参数 $\varepsilon$ 保留在求解目标中。

位置变换矩阵 $\bm T_k\in\mathbb R^{N\times N}$ 按行对应源粒子 $i$，按列对应目标粒子 $j$，定义为
\begin{equation}\label{cw:eq:particle-ot-conditional}
\begin{aligned}
&(\bm T_k)_{ij}:=N(\bm\Pi_k)_{ij},\;\bm T_k=N\bm\Pi_k,\\
&\bm T_k^{\mathrm T}\mathbf1_N=\mathbf1_N,\;\bm T_k\mathbf1_N=N\bm w_k^+.
\end{aligned}
\end{equation}
每列的变换系数非负且和为 $1$, 据此取源粒子的加权平均, 得到新粒子：
\begin{equation}\label{cw:eq:particle-ot-transform}
\begin{aligned}
&x_k^{+,(j)}:=\sum_{i=1}^N(\bm T_k)_{ij}x_k^{-,(i)}=N\sum_{i=1}^N(\bm\Pi_k)_{ij}x_k^{-,(i)},\\
&X_k^+=X_k^-\bm T_k=NX_k^-\bm\Pi_k.
\end{aligned}
\end{equation}

% !TeX root = ../../main.tex
\section{压缩与还原流形拟合}\label{cw:sec:ml-compress}

当观测取值于高维向量空间 $\mathcal S_X=\mathbb R^{n_D}$ 时，直接在环境空间中实施似然加权或离散耦合输运会遭遇严重的维数灾难。实际物理或视觉信号往往并不充满整个环境空间，而是集中于某个自由度仅为 $n_d$（$n_d<n_D$）的低维子集 $\mathcal M\subseteq\mathcal S_X$ 附近。本节从经验模内积下的重构误差出发，给出压缩算子 $\mathscr T_{\Rightarrow}:\mathcal S_X\to\mathcal S_Z$（$\mathcal S_Z=\mathbb R^{n_d}$）与还原算子 $\mathscr T_{\Leftarrow}:\mathcal S_Z\to\mathcal S_X$ 的流形拟合框架，为第~\ref{cw:sec:algorithms}~章的潜空间输运滤波提供低维状态坐标与经验推前先验。

设离线观测样本矩阵为 $x_{1..M}=(x_1,\dots,x_M)\in\mathbb R^{n_D\times M}$。在样本行排列空间上定义经验模内积 $\langle\bm U,\bm V\rangle_{\hat{\mathbb P}_M}:=\frac1M\bm U\bm V^{\mathrm T}\in\mathbb R^{n_D\times n_D}$ 及其诱导范数 $\|\bm U\|_{\hat{\mathbb P}_M}^2:=\operatorname{tr}\langle\bm U,\bm U\rangle_{\hat{\mathbb P}_M}$。对任意重构映射 $\mathscr R:\mathcal S_X\to\mathcal M$，记余映射为 $\mathscr R^{\perp}:=\mathscr I-\mathscr R$，则样本平均重构平方误差可统一表示为
\begin{equation}\label{cw:eq:manifold-emp-norm}
\frac1M\sum_{i=1}^M\bigl\|x_i-\mathscr R(x_i)\bigr\|^2
=\operatorname{tr}\bigl\langle\mathscr R^{\perp}(x_{1..M}),\mathscr R^{\perp}(x_{1..M})\bigr\rangle_{\hat{\mathbb P}_M}
=\bigl\|\mathscr R^{\perp}(x_{1..M})\bigr\|_{\hat{\mathbb P}_M}^2.
\end{equation}
若不对 $\mathscr R$ 施加维数瓶颈而直接极小化式~\eqref{cw:eq:manifold-emp-norm}，最优解将平凡退化为恒等算子 $\mathscr I$。引入以 $\mathcal S_Z$ 为中介的压缩—还原复合结构 $\mathscr R=\mathscr T_{\Leftarrow}\circ\mathscr T_{\Rightarrow}$，即可将重构像集限制在至多 $n_d$ 个自由度的低维几何结构上。

\subsection{平直线性子空间的伴随投影与谱截断}

先考察 $\mathcal M$ 为 $n_d$ 维平直真子空间 $W<\mathcal S_X$ 的基础情形。此时压缩与还原算子均为线性映射，标准基下的表示矩阵分别记为 $\bm T_{\Rightarrow}\in\mathbb R^{n_d\times n_D}$ 与满列秩矩阵 $\bm T_{\Leftarrow}\in\mathbb R^{n_D\times n_d}$，像空间 $W:=\operatorname{im}(\mathscr T_{\Leftarrow})<\mathcal S_X$ 自然满足 $\operatorname{rank}(\bm T_{\Leftarrow}\bm T_{\Rightarrow})\le n_d$。记向 $W$ 的正交投影算子为 $\mathscr P_W$（矩阵表示为 $\bm P_W$），余投影为 $\mathscr P_W^{\perp}=\mathscr I-\mathscr P_W$（矩阵表示为 $\bm P_W^{\perp}=\bm I_{n_D}-\bm P_W$）。利用 $\operatorname{im}(\bm P_W^{\perp})\perp W$ 及迹的轮换性，重构误差可作正交分解：
\begin{equation}\label{cw:eq:manifold-orth-decomp}
\bigl\|(\bm I_{n_D}-\bm T_{\Leftarrow}\bm T_{\Rightarrow})x_{1..M}\bigr\|_{\hat{\mathbb P}_M}^2
=\operatorname{tr}\bigl(\bm P_W^{\perp}\langle x_{1..M},x_{1..M}\rangle_{\hat{\mathbb P}_M}\bigr)
+\bigl\|(\bm P_W-\bm T_{\Leftarrow}\bm T_{\Rightarrow})x_{1..M}\bigr\|_{\hat{\mathbb P}_M}^2.
\end{equation}
式~\eqref{cw:eq:manifold-orth-decomp} 右端首项仅依赖像子空间 $W$ 的选取，次项则度量复合映射偏离正交投影的程度。由此可直接导出线性压缩与子空间选取的闭式最优解，其完整推导见附录~\ref{cw:app:manifold}。

\begin{proposition}[线性压缩还原的伴随闭式解与谱截断]\label{cw:prop:linear-manifold}
设 $\mathscr T_{\Leftarrow}:\mathcal S_Z\to\mathcal S_X$ 为单射线性算子，$W=\operatorname{im}(\mathscr T_{\Leftarrow})$，$\mathscr T_{\Leftarrow}^{\star}:\mathcal S_X\to\mathcal S_Z$ 为其内积伴随算子。记样本二阶矩矩阵 $\langle x_{1..M},x_{1..M}\rangle_{\hat{\mathbb P}_M}$ 的特征值按降序排列为 $\lambda_1\ge\dots\ge\lambda_{n_D}\ge0$，相应规范正交特征向量族为 $u_{1..n_D}$。则：
\begin{enumerate}[label=(\arabic*),leftmargin=2.2em,topsep=2pt,itemsep=2pt]
\item 固定 $\mathscr T_{\Leftarrow}$ 时，最优压缩算子为 $\mathscr T_{\Rightarrow}^{\optstar}=(\mathscr T_{\Leftarrow}^{\star}\mathscr T_{\Leftarrow})^{-1}\mathscr T_{\Leftarrow}^{\star}$（标准实内积下矩阵形式为 $\bm T_{\Rightarrow}^{\optstar}=(\bm T_{\Leftarrow}^{\mathrm T}\bm T_{\Leftarrow})^{-1}\bm T_{\Leftarrow}^{\mathrm T}$），此时复合算子恰好还原正交投影 $\mathscr T_{\Leftarrow}\mathscr T_{\Rightarrow}^{\optstar}=\mathscr P_W$，使式~\eqref{cw:eq:manifold-orth-decomp} 次项归零；
\item 在全体 $n_d$ 维真子空间 $W<\mathcal S_X$ 中，前 $n_d$ 个主特征向量张成的子空间 $W^{\optstar}=\operatorname{span}\{u_1,\dots,u_{n_d}\}$（对应 $\bm P_{W^{\optstar}}=\sum_{j=1}^{n_d}\bm P_{u_j}$，其中 $\bm P_{u_j}:=u_ju_j^{\mathrm T}$）使重构误差达到全局极小值 $\sum_{j=n_d+1}^{n_D}\lambda_j$，且当 $\lambda_{n_d}>\lambda_{n_d+1}$ 时最优子空间唯一。
\end{enumerate}
\end{proposition}

\subsection{非线性参数化流形与推前先验}

当高维信号的内在几何呈显著弯曲形态（如手写数字图像随笔画形变构成的非线性曲面）时，平直子空间 $W$ 需要较高的维数才能压低残差特征值之和。为此，将压缩与还原算子推广为由参数 $\bm\theta\in\mathcal S_{\bm\Theta}$ 索引的神经网络映射族 $\mathscr T_{\Rightarrow}^{\bm\theta}:\mathcal S_X\to\mathcal S_Z$ 与 $\mathscr T_{\Leftarrow}^{\bm\theta}:\mathcal S_Z\to\mathcal S_X$，端到端重构映射记为 $\mathscr R^{\bm\theta}:=\mathscr T_{\Leftarrow}^{\bm\theta}\circ\mathscr T_{\Rightarrow}^{\bm\theta}$。还原算子的像集 $\mathcal M_{\bm\theta}:=\operatorname{im}(\mathscr T_{\Leftarrow}^{\bm\theta})$ 在环境空间 $\mathcal S_X$ 中给出由潜坐标 $z\in\mathcal S_Z$ 参数化的低维曲面；在任意潜坐标 $z_0\in\mathcal S_Z$ 的邻域内作一阶局部线性展开：
\begin{equation}\label{cw:eq:manifold-tangent}
\mathscr T_{\Leftarrow}^{\bm\theta}(z_0+\delta z)
\approx\mathscr T_{\Leftarrow}^{\bm\theta}(z_0)
+\frac{\mathrm d\mathscr T_{\Leftarrow}^{\bm\theta}}{\mathrm dz^{\mathrm T}}(z_0)\,\delta z,
\end{equation}
导数矩阵 $\frac{\mathrm d\mathscr T_{\Leftarrow}^{\bm\theta}}{\mathrm dz^{\mathrm T}}(z_0)\in\mathbb R^{n_D\times n_d}$ 的列空间即张成流形在该点处的局部切空间，当映射退化为常矩阵 $\bm T_{\Leftarrow}$ 时即回到命题~\ref{cw:prop:linear-manifold} 的平直情形。一般非线性参数下问题无闭式解，通过随机梯度优化联合极小化经验重构误差确定参数：
\begin{equation}\label{cw:eq:ml-reconstruction}
\bm\theta^{\optstar}\in\operatorname*{argmin}_{\bm\theta\in\mathcal S_{\bm\Theta}}\
\bigl\|(\mathscr I-\mathscr R^{\bm\theta})(x_{1..M})\bigr\|_{\hat{\mathbb P}_M}^2
=\operatorname*{argmin}_{\bm\theta\in\mathcal S_{\bm\Theta}}\
\frac1M\sum_{i=1}^M\bigl\|x_i-\mathscr T_{\Leftarrow}^{\bm\theta}\bigl(\mathscr T_{\Rightarrow}^{\bm\theta}(x_i)\bigr)\bigr\|^2.
\end{equation}
为避免过拟合或约束潜空间几何尺度，亦可在式~\eqref{cw:eq:ml-reconstruction} 中附加参数范数等正则项。下文固定训练所得参数 $\bm\theta^{\optstar}$，在不致混淆处将 $\mathscr T_{\Rightarrow}^{\bm\theta^{\optstar}},\mathscr T_{\Leftarrow}^{\bm\theta^{\optstar}}$ 简记为 $\mathscr T_{\Rightarrow},\mathscr T_{\Leftarrow}$。

确定压缩算子后，直接取训练样本经验测度 $\hat{\mathbb P}_{M;X}$ 在 $\mathscr T_{\Rightarrow}$ 下的推前测度作为潜状态的先验分布：
\begin{equation}\label{cw:eq:ml-prior}
\mathbb P_Z:=(\mathscr T_{\Rightarrow})_*\hat{\mathbb P}_{M;X}
=\frac1M\sum_{i=1}^M\delta_{\mathscr T_{\Rightarrow}(x_i)}.
\end{equation}
这一构造使潜空间初始粒子直接取自真实样本的压缩坐标，忠实保留了数据流形上的多峰聚集结构，无须强行假定潜变量服从标准正态分布。

% !TeX root = ../../main.tex
\section{条件输运映射与潜空间后验表示}\label{cw:sec:ml-transport-theory}

第~\ref{cw:sec:ot}~节的离散耦合变换在给定预测支撑点上重分配质量，仍以重要性权重作为源边际输入。本节从连续映射视角出发，建立不经由似然权重的条件最优输运极大极小对偶表述，并在第~\ref{cw:sec:ml-compress}~节的潜空间上给出局部序贯观测模型、最优后验分类准则及静态潜状态后验的一步累积表示，作为第~\ref{cw:sec:ml-updates}~节算法构造与机制分析的理论基础。

\subsection{条件 Monge 推前约束与极大极小对偶表述}\label{cw:sec:ml-minimax}

考虑潜状态与观测二元随机变量 $(Z,Y)\sim\mathbb P_{(Z,Y)}$，另取与 $(Z,Y)$ 独立的参考先验随机变量 $\bar Z\sim\mathbb P_Z$，使独立耦合满足 $(\bar Z,Y)\sim\mathbb P_Z\otimes\mathbb P_Y$。条件输运方法旨在寻找可测映射 $T:\mathcal S_Z\times\mathcal S_Y\to\mathcal S_Z$，使对 $\mathbb P_Y$-几乎处处的观测实现 $y\in\mathcal S_Y$ 均成立条件推前关系 $T(\cdot,y)_*\mathbb P_Z=\mathbb P_{Z\mid Y=y}$\cite{aljarrah2023otpf,aljarrah2024brenier}。引入保持观测坐标不变的联合提升映射 $S_T(z,y):=(T(z,y),y)$，结合条件核乘积分解，逐点条件推前关系等价于联合测度推前等式：
\begin{equation}\label{cw:eq:ml-joint-constraint}
(S_T)_*(\mathbb P_Z\otimes\mathbb P_Y)=\mathbb P_{(Z,Y)}.
\end{equation}
式~\eqref{cw:eq:ml-joint-constraint} 表明，将独立先验样本 $\bar Z$ 连同实际观测 $Y$ 送入映射后，所得二元组 $(T(\bar Z,Y),Y)$ 与真实状态—观测对 $(Z,Y)$ 同分布。满足式~\eqref{cw:eq:ml-joint-constraint} 的可测映射一般并不唯一；与式~\eqref{cw:eq:monge} 的二次代价准则一致，在可行映射族中选取平均平方位移最小者：
\begin{equation}\label{cw:eq:ml-constrained}
\min_{T}\ \mathbb E_{(\bar Z,Y)\sim\mathbb P_Z\otimes\mathbb P_Y}\Bigl[\tfrac12\bigl\|T(\bar Z,Y)-\bar Z\bigr\|^2\Bigr]
\qquad\text{s.t.}\qquad (S_T)_*(\mathbb P_Z\otimes\mathbb P_Y)=\mathbb P_{(Z,Y)}.
\end{equation}

与第~\ref{cw:sec:sinkhorn}~节对离散边缘约束引入乘子的思路平行，利用连续有界检验函数 $f:\mathcal S_Z\times\mathcal S_Y\to\mathbb R$ 刻画测度等式~\eqref{cw:eq:ml-joint-constraint}。由推前与拉回的对偶配对关系
\begin{equation}
\langle f,(S_T)_*(\mathbb P_Z\otimes\mathbb P_Y)\rangle_{\mathcal S_Z\times\mathcal S_Y}
=\langle S_T^*f,\mathbb P_Z\otimes\mathbb P_Y\rangle_{\mathcal S_Z\times\mathcal S_Y},
\end{equation}
将检验函数 $f$ 作为约束乘子并入目标，定义对偶鞍点泛函：
\begin{equation}\label{cw:eq:ml-objective}
\begin{aligned}
J(f,T):={}&\langle f,\mathbb P_{(Z,Y)}\rangle_{\mathcal S_Z\times\mathcal S_Y}
-\langle S_T^*f,\mathbb P_Z\otimes\mathbb P_Y\rangle_{\mathcal S_Z\times\mathcal S_Y}\\
&+\mathbb E_{(\bar Z,Y)\sim\mathbb P_Z\otimes\mathbb P_Y}\Bigl[\tfrac12\bigl\|T(\bar Z,Y)-\bar Z\bigr\|^2\Bigr].
\end{aligned}
\end{equation}
当推前约束式~\eqref{cw:eq:ml-joint-constraint} 不成立时，$\sup_f J(f,T)=+\infty$，故原约束问题~\eqref{cw:eq:ml-constrained} 等价于 $\min_T\sup_f J(f,T)$。交换极值次序后得到无约束极大极小鞍点问题：
\begin{equation}\label{cw:eq:ml-minimax}
(f^{\optstar},T^{\optstar})\in\operatorname*{argmax}_{f}\ \operatorname*{argmin}_{T}\ J(f,T).
\end{equation}
在式~\eqref{cw:eq:ml-objective} 中，前两项之差度量推前联合分布与真实联合分布在检验函数 $f$ 下的期望偏差，外层极大化寻找最具区分度的判别函数，内层极小化驱动推前样本匹配真实后验；末项则以二次位移代价固定最优映射的几何选择。关于第一变元的凸势约束与极值交换合法性的严格条件参见文献\cite{aljarrah2024brenier}。值得注意的是，泛函 $J(f,T)$ 仅通过联合样本 $(Z,Y)$ 与独立样本 $\bar Z$ 的期望定义，不涉及观测条件密度的显式解析求值。

\subsection{潜空间局部序贯观测模型与最优识别准则}\label{cw:sec:ml-problem}

将第~\ref{cw:sec:ml-compress}~节的流形表示用于局部序贯观测下的模式识别。设完整高维对象 $X\in\mathcal S_X=\mathbb R^{n_D}$ 及其离散类别 $C\in\{0,\dots,9\}$ 由潜空间坐标 $Z\in\mathcal S_Z=\mathbb R^{n_d}$（$n_d<n_D$）经还原算子 $\mathscr T_{\Leftarrow}:\mathcal S_Z\to\mathcal S_X$ 与读出映射 $\ell:\mathcal S_Z\to\{0,\dots,9\}$ 确定：
\begin{equation}\label{cw:eq:ml-latent}
X=\mathscr T_{\Leftarrow}(Z),\qquad C=\ell(Z),\qquad Z\sim\mathbb P_Z,
\end{equation}
其中先验 $\mathbb P_Z$ 取式~\eqref{cw:eq:ml-prior} 的经验推前测度。在序贯观测过程中，底层对象保持静止，时刻 $k\in\{1,\dots,K\}$ 仅经由局部选择矩阵 $\bm H_k\in\{0,1\}^{r^2\times n_D}$ 观测第 $k$ 个 $r\times r$ 窗口的含噪响应：
\begin{equation}\label{cw:eq:ml-model}
\begin{aligned}
&Z_k=Z_{k-1},\qquad Z_0\sim\mathbb P_Z,\\
&Y_k=\bm H_k\mathscr T_{\Leftarrow}(Z_k)+N_{y,k},\qquad N_{y,k}\sim\mathcal N(\bm 0,\sigma^2\bm I_{r^2}),
\end{aligned}
\end{equation}
各时刻观测噪声 $\{N_{y,k}\}$ 相互独立且与 $Z_0$ 独立。式~\eqref{cw:eq:ml-model} 是第~\ref{cw:sec:model}~节一般非线性滤波模型在潜状态静止（过程噪声恒为零）下的具体实例。

\begin{proposition}[序贯识别的最优后验准则]\label{cw:lem:ml-bayes-rule}
在生成关系~\eqref{cw:eq:ml-latent} 与状态空间模型~\eqref{cw:eq:ml-model} 下，时刻 $k$ 使错判概率 $\mathbb P\bigl(\hat c_k(Y_{1:k})\neq C\bigr)$ 最小的可测分类准则满足
\begin{equation}\label{cw:eq:ml-bayes-rule}
\hat c_k^{\optstar}(y_{1:k})\in\operatorname*{argmax}_{c\in\{0,\dots,9\}}\
\mathbb P_{Z_k\mid Y_{1:k}=y_{1:k}}\bigl(\ell^{-1}(c)\bigr).
\end{equation}
\end{proposition}
\begin{proof}
对任意可测分类器 $\hat c_k$，正确分类概率满足全期望分解
$\mathbb P\bigl(\hat c_k(Y_{1:k})=C\bigr)=\mathbb E_{Y_{1:k}\sim\mathbb P_{Y_{1:k}}}\bigl[\mathbb P\bigl(C=\hat c_k(Y_{1:k})\mid Y_{1:k}\bigr)\bigr]$，
对每个实现 $y_{1:k}$ 逐点最大化条件后验概率即可最小化总错判率。由 $C=\ell(Z_0)=\ell(Z_k)$，有 $\mathbb P(C=c\mid Y_{1:k}=y_{1:k})=\mathbb P_{Z_k\mid Y_{1:k}=y_{1:k}}(\ell^{-1}(c))$。
\end{proof}

当滤波后验由等权经验测度 $\frac1N\sum_{i=1}^N\delta_{z_k^{+,(i)}}$ 近似时，将经验测度代入式~\eqref{cw:eq:ml-bayes-rule} 即得粒子读出众数判决：
\begin{equation}\label{cw:eq:ml-vote}
\hat c_k=\operatorname*{argmax}_{c\in\{0,\dots,9\}}\ \frac1N\sum_{i=1}^N\mathbb I_{\{c\}}\bigl(\ell(z_k^{+,(i)})\bigr).
\end{equation}

\subsection{静止潜状态下滤波后验的一步累积表示}\label{cw:sec:ml-static-theory}

一般动态系统中，第 $k$ 步预测分布 $\mathbb P_{Z_k\mid Y_{1:k-1}=y_{1:k-1}}$ 依赖历史观测路径，使式~\eqref{cw:eq:ml-minimax} 中的源分布随步改变。在式~\eqref{cw:eq:ml-model} 的静态演化下，多步递推后验可精确化归为固定先验上的一步条件更新。

\begin{proposition}[静止状态下滤波后验的一步表示]\label{cw:lem:ml-static}
设潜状态满足 $Z_k=Z_0\sim\mathbb P_Z$，且给定潜状态时各时刻观测 $\{Y_j\}_{j=1}^k$ 条件独立，条件密度记为 $p_{Y_j\mid Z_j}(y_j\mid z)$。当分母积分有限且严格为正时，对任意时刻 $k\ge1$ 有
\begin{equation}\label{cw:eq:ml-static}
\frac{\mathrm d\mathbb P_{Z_k\mid Y_{1:k}=y_{1:k}}}{\mathrm d\mathbb P_Z}(z)
=\frac{\prod_{j=1}^k p_{Y_j\mid Z_j}(y_j\mid z)}{\int_{\mathcal S_Z}\prod_{j=1}^k p_{Y_j\mid Z_j}(y_j\mid\tilde z)\,\mathrm d\mathbb P_Z(\tilde z)},
\end{equation}
即第 $k$ 步递推滤波后验等于以累积观测 $y_{1:k}$ 对初始先验 $\mathbb P_Z$ 作单次更新所得的后验分布。
\end{proposition}
\begin{proof}
对 $k$ 归纳。$k=1$ 时由定理~\ref{cw:thm:bayes} 成立。设结论对 $k-1$ 成立；因 $Z_k=Z_{k-1}$，预测分布满足 $\mathbb P_{Z_k\mid Y_{1:k-1}=y_{1:k-1}}=\mathbb P_{Z_{k-1}\mid Y_{1:k-1}=y_{1:k-1}}$。在时刻 $k$ 应用条件独立性与密度导数链式法则，将该分布关于 $\mathbb P_Z$ 的密度导数乘以 $p_{Y_k\mid Z_k}(y_k\mid z)$ 并归一化，即得式~\eqref{cw:eq:ml-static}。
\end{proof}

记 $y\in\mathbb R^{n_D}$ 为叠加全图观测噪声后的像素向量，$\bm m_k\in\{0,1\}^{n_D}$ 为标记前 $k$ 个窗口覆盖像素位置的二值掩码向量，则变长累积观测 $y_{1:k}$ 可统一嵌入定长空间表示为 $(\bm m_k\odot y,\bm m_k)$（未揭示区域置零，并由 $\bm m_k$ 区分未观测像素与零值背景）。结合命题~\ref{cw:lem:ml-static} 与式~\eqref{cw:eq:ml-joint-constraint}，只需构造单一的摊销条件输运映射 $T:\mathcal S_Z\times\mathbb R^{n_D}\times\{0,1\}^{n_D}\to\mathcal S_Z$，使其对一切时刻 $k\in\{1,\dots,K\}$ 满足一步推前关系
\begin{equation}\label{cw:eq:ml-amortized-pushforward}
T(\cdot,\bm m_k\odot y,\bm m_k)_*\mathbb P_Z=\mathbb P_{Z_k\mid Y_{1:k}=y_{1:k}}.
\end{equation}


\section{本章小结}

本章梳理了全文所需的数学基础。在一般非线性系统中，精确的 Bayes 滤波没有解析公式，粒子滤波于是改用有限个带权重的样本去近似后验分布。问题在于，一旦样本权重严重失衡，传统多项重采样只会“留大弃小、原位复制”，无法把粒子挪到新的位置上；离散最优输运则换了一种思路——先算出预测粒子与目标后验之间最省力的配对方案（耦合矩阵），再用加权平均把粒子平滑地搬移过去，从而在不依赖随机复制的前提下完成重采样。

当观测维数升高（如整幅图像）时，仅靠平移现有粒子仍不够用，因为高维空间里的似然权重本身就会迅速退化失效。为此，本章引入压缩与还原流形拟合：在线性情形下它等价于向主特征子空间作正交投影，在非线性情形下则用神经网络拟合高维对象的低维流形以提取紧凑潜坐标，并直接由训练数据给出潜空间的先验分布。在此基础上，条件输运映射利用极大极小对抗目标直接把先验样本推向观测条件下的后验位置，彻底绕开了权重计算；特别是当待识别的潜状态本身不随时间改变时，任意时刻的后验只取决于截至目前已看到的累积观测，因而无论经历多少步观测，都可以从初始粒子一步映射到位。这一性质直接构成了下一章离线摊销滤波算法的理论支点。
